\section{Method Chaining}
Some APIs are designed to allow \bdq{method chaining}\index{method chaining}, sometimes also referred as \bdq{Call Chaining}\index{call chaining|see {method chaining}}; you may find this in particular for implementations of the \bdq{builder} pattern\index{builder pattern}, but it appears elsewhere, too. A prominent example is the class \lstinline|java.lang.StringBuilder|\index{java.lang!StringBuilder} \autocite{ORACLE_DOC_STRINGBUILDER_CLASS} (obviously not a builder, despite its name), a real example for a builder class would be the class \lstinline|DateTimeFormatterBuilder|\index{java.time.format!DateTimeFormatterBuilder} \autocite{ORACLE_DOC_DATETIMEFORMATTERBUILDER_CLASS} from the package \lstinline|java.time.format|\index{java.time.format}.

Usually, a class (or interface) will designed for method chaining\index{method chaining} when its methods return \lstinline|this| (the containing object instance), but other designs are possible that also allow method chaining\index{method chaining}, as shown in the code snippet below. It shows the use of method chaining\index{method chaining} on streams and on an instance of \lstinline|Optional|\index{java.util!Optional}.

\keeplisting{6}
\index{java.util!Map}\index{java.util!Map!Entry}\index{java.util!Optional}\index{java.util!Stream}
\begin{lstlisting}
return strings.stream()
    .collect( groupingBy( s -> s, counting() ) ) // returns java.util.Map
    .entrySet() 
    .stream()
    .max( Map.Entry.comparingByValue() ) // returns java.util.Optional
    .map( Map.Entry::getKey );
\end{lstlisting}

\keeplisting{3}
Method chaining allows you to write
\index{java.io!Writer}\index{java.io!Writer!append()}
\begin{lstlisting}
writer.append( "Caption\t: " )
    .append( value );
\end{lstlisting}

\keeplisting{3}
instead of
\index{java.io!Writer}\index{java.io!Writer!append()}
\begin{lstlisting}
writer.append( "Caption\t: " );
writer.append( value );
\end{lstlisting}

If the \bdq{chain} does not fit completely into a single line, each method call has to be written into a single line of its own.

\keeplisting{7}
In fact, it is recommended to do this all the time:
\index{java.io!Writer}\index{java.io!Writer!append()}
\begin{lstlisting}
// Acceptable
writer.append( "Caption\t: " ).append( value );

// RECOMMENDED
writer.append( "Caption\t: " )
    .append( value );
\end{lstlisting}

The dot has to be written in front of the method's name and the methods will be indented regularly.

\keeplisting{8}
Streams, that had been introduced with Java~8, will be formatted in the same way:
\index{java.util!Stream}
\begin{lstlisting}
final var names = customers.stream()
    .filter( c -> c.getCountry().equals( GERMANY )
    .filter( c -> c.getTurnover() > limit )
    .map( Customer::getName )
    .sorted()
    .toArray( String []::new );
\end{lstlisting}


\subsubsection{Principles}
\begin{enumerate}[label=P\arabic*.]
\item{Each method call after the first should be placed on a line of it own.

The dot prepends the method's name.}
\item{The method calls are indented one level relative to the first line of the call chain.} 
\end{enumerate}

%==============================================================================
% End of file!!
%==============================================================================
